% Pursuing Stacks (cotes 134-2 to 134-8): archive cote 134-3, site page 85 (PDF page 86), typescript page 85.
% typescript. Draft compared: ps2.tex:4434-4498.
% AI-transcribed by Claude Opus 5.5 (claude-opus-5-5, adaptive thinking, effort high), Message Batches API, batch ps1 (msgbatch_01Rv5hCaNTtUt4JHx1yqoHPL), 2026-11-12; not yet checked by eye.
\scanpage{134-3/85}
\mspage{85}
$\hand{\Delta}^{f}$ of standard ordered simplices with only ``face-like\unsure{''|*} maps between them\hand{,}
namely strictly increasing ones (that is, ruling out degeneracies) is
a weak test category, and not a test category. On the other hand, for
any test category A, we can construct lots of test categories which are
not strict, namely all categories $A_{/F}$ where F is any \inserted[typed]{aspheric\hand{,}} object in $A^{\wedge}$ which
admits two ``non-empty'' subobjects whose intersection is ``empty'' - take
for instance for F any homotopy interval. In the case of $A = \hand{\Delta}$ , we may
take for F any object $\hand{\Delta}_{n}$ $(n \hand{\geq} 1)$ in A, except \inserted{just} the final object $\hand{\Delta}_{o}$ .
\deleted{\uline{Remark} Coming back to the notion of a weak test category, still another}
\deleted{equivalent formulation occured to me while writing up these reflections,}
\deleted{namely}

\deleted{(i\unsure{x}) The functor $j_{A}$: (Cat) \quad $A^{\wedge}$ transforms weak equivalences into}
\deleted{weak equivalences.}

\deleted{This condition is trivially implied by (i), let's prove it implies}
\deleted{(v) namely that for C in (Cat), C with final element, $j_{A}(C)$ is aspheric.}
\deleted{As C \quad e is a weak equivalence so is $j_{A}(C)$ \quad $j_{A}(e) = e_{A^{\wedge}}$ by assumption}
\deleted{thus all we have to check is that $e_{A^{\wedge}}$ is aspheric.}

d) \uline{Test functors}. Let A be a weak test category, and let
\[ (9) \qquad i:\ A \hand{\longrightarrow} (\mathrm{Cat}) \]
be a functor, giving rise to a functor
\[ (10) \qquad i^{*} :\ (\mathrm{Cat}) \hand{\longrightarrow} A^{\wedge} , \quad C \hand{\longmapsto} i^{*}(C) = (a \hand{\longmapsto} \mathrm{Hom}(i(a),C)\hand{)} \]
We'll say that i is a \inserted{weak} \uline{test functor} if $i^{*}$ is \xout{\unsure{not}} a morphism of modelize\unsure{rs}
i.e. model-preserving, namely
\[ \hand{\Big\{}\begin{array}{l} W_{(\mathrm{Cat})} = (i^{*})^{-1}(W_{A^{\wedge}}) \quad \text{and} \\ W_{(\mathrm{Cat})}{}^{-1}\,(\mathrm{Cat}) = (\mathrm{Hot}) \hand{\longrightarrow} W_{A^{\wedge}}{}^{-1}\,A^{\wedge} \text{ an equivalence} \quad . \end{array} \]
As by assumption we know already that (1)
\[ i_{A} :\ A^{\wedge} \hand{\longrightarrow} (\mathrm{Cat}) \quad , \quad F \hand{\longmapsto} A_{/F} \]
is model-preserving, this implies that i is a \inserted{weak} test functor iff the
compositum
\[ (11) \qquad i_{A}i^{*} :\ (\mathrm{Cat}) \hand{\longrightarrow} (\mathrm{Cat}) \quad , \quad C \hand{\longmapsto} A_{/i^{\hand{*}}(C)} \overset{\mathrm{def}}{=} A_{/C} \]
is model-preserving, i.e. \inserted{(essentially)} induces an \uline{autoequivalence} of (Hot).

The basic example is to take $i = i_{A}$ , hence $i^{*} = j_{A}$ (2) , the
\xout{very} fact that A is a weak test category can be just translated \inserted{(p.80, (vii))} by saying
that this functor (11) = (6) \inserted[typed]{transforms weak equivalences into weak} \xout{is isomorphic to the identity functor,} \xout{\illegible{1}}
equivalent\xout{\unsure{es}}, \xout{transform} and the localized functor (Hot)$\hand{\longrightarrow}$(Hot) is an
equivalence (or, what is enough, transforms final object into final
object of (Hot)). But as we saw\inserted{before}\hand{,} this \inserted{weak} test functor gives rise to
categories $i(a) = A_{/a}$ which are prohibitively large, and one generally
prefers working with \inserted{weak} test functors more appropriate for computations -
